% Converted from markdown/graph_organization.md

\section{Update-Friendly Heterogeneous Graph Organization}


Topological regionality offers an opportunity to limit adjacency movement and GPU index maintenance to the sources whose outgoing adjacency changes in a batch. However, insertion operations in the PMA structure may trigger rebalancing, causing the adjacency list to be migrated to a region other than that of the source vertex, thereby increasing the volume of data movement and the overhead of index maintenance at other locations. We organize maintenance around independently allocated source blocks. Section 4.1 explains how this layout prevents updates from relocating unrelated lists. Section 4.2 coordinates updates to outgoing and incoming adjacency through a common batch plan. Section 4.3 uses the resulting changed-source set to update only the required GPU adjacency indices while preserving consistent graph access.

\subsection{Source-Isolated Graph Layout}


The graph resides in host memory. For each source, the GPU maintains an adjacency descriptor, an index record that identifies where its neighbor list is stored and how many neighbors it contains. Although graph updates are regional, modifying a few vertices can relocate the adjacency lists of many unchanged vertices, requiring corresponding updates to their GPU indices and thereby amplifying the data update overhead. To confine this maintenance to updated sources, we allocate and grow each source's adjacency independently.

Each source $u$ stores its neighbors in an independently allocated, contiguous block of mapped host memory, with spare capacity for insertions. For insertions, the CPU appends new neighbors to $u$'s existing block if its capacity is sufficient. Otherwise, it allocates a larger block, copies $u$'s existing adjacency into it, and appends the new neighbors. Expansion uses a separately allocated block, leaving the locations of other adjacency lists unchanged. Fig.~\ref{fig:2}(a) contrasts this with a shared PMA: even a nearby shift before rebalancing can change the start index of an unchanged source $v$, whereas expanding $u$'s independent block leaves $v$'s block and descriptor valid. For deletions, the CPU removes the specified edges that exist in the graph and compacts the remaining neighbors within the existing block. This preserves the contiguous neighbor access of the original layout while restricting relocation to the updated source.


Independent block allocation eliminates cross-source redistribution, preserving the adjacency locations and descriptors of unchanged sources. This confines relocation costs to updated lists and enables sparse GPU publication (Section~4.3). Expansion and deletion may still copy or scan an entire updated list, while spare capacity and delayed block reclamation incur memory overhead.

\subsection{Unified Change View}


A topology view is a representation of graph in a particular data structure or on a particular device, organized for the access it supports. Our system maintains three such views. Outgoing adjacency stores each source's neighbors in host memory for traversing its outgoing edges. The reverse index stores the sources of incoming edges for each destination, allowing deletion recovery to find the remaining predecessors that may restore its distance. GPU adjacency descriptors store each source's neighbor-list location and degree, allowing GPU computation to locate and read its outgoing neighbors. Each topology update must be reflected in all three views. We coordinate their maintenance through a unified change view, consisting of effective edge-change records and changed-source lists produced during outgoing-adjacency preparation for each update phase.

The CPU groups update requests by source so that adjacency maintenance operates on a group of edges at a time. For each source, it collects all deletion requests and records how many existing edges to each destination should be removed. It then removes these edges together by compacting the neighbor list once, avoiding repeated movement of the same entries when edges are deleted one at a time. It also determines the capacity required for the source's insertions as a group. If the current block is too small, the CPU allocates a larger block and copies the existing neighbors once before appending the new edges. Thus, grouping reduces repeated movement within an updated adjacency list, complementing the isolation of movement across sources in Section 4.1. Independent source blocks allow the compaction and append operations for different sources to proceed in parallel.

Grouping updates by source also provides the information needed to maintain the other topology views. During preparation, the CPU records each source--destination pair and the number of edges actually removed or added, together with the sources whose adjacency changes. Unsuccessful deletion requests produce no change. The reverse index uses the edge records to update incoming adjacency, while GPU descriptor and cache updates use the source lists to locate the outgoing lists that have changed. These structures therefore reuse the results of adjacency preparation instead of independently determining the effects of the same requests.

For example, suppose source $u$ has neighbors $[a,a,b]$ and receives three requests to delete $(u,a)$ and one to insert $(u,c)$. Only two copies of $(u,a)$ exist, so deletion compacts the list to $[b]$ and produces the record $(u,a,-2)$. The reverse index uses this record to remove the two incoming edges from $u$ at $a$. Insertion later produces $[b,c]$ and the record $(u,c,+1)$, which adds an incoming edge from $u$ at $c$. Both stages identify $u$ as changed. The system then updates its GPU descriptor and, if present, its cached adjacency, as illustrated in Fig.~\ref{fig:2}(b). The third deletion request causes no additional index update.


\begin{figure*}[t]
\centering
\fbox{\parbox[c][0.85in][c]{0.9\textwidth}{\centering Figure artwork to be added}}
\caption{Shared-array movement and source-isolated maintenance. (a) In the illustrative PMA, insertion first shifts an unchanged source's adjacency, a later leaf overflow rebalances the eight-slot range, and root overflow expands it to 16 slots. The corresponding source-local updates leave the unchanged source's address stable and expand only the updated source's block. (b) Effective deletion and insertion records update the reverse overlay; their changed-source union selects one final GPU descriptor and resident-cache repair after deletion recovery. The two panels use separate examples.}
\label{fig:2}
\end{figure*}


\subsection{Maintaining Reverse and GPU Views}


\textbf{Reverse index maintenance.} Maintaining incoming adjacency directly from the input requests would cause other views to repeat checks already performed when processing outgoing adjacency. To avoid this work, the reverse index consumes the actual edge changes, groups them by destination, and combines changes for the same source--destination pair. Each destination stores a sorted base list of incoming sources and a separate list of delta records. A delta record contains a source vertex and the net number of edges added or removed relative to the base list: a positive value indicates additions, and a negative value indicates deletions. New changes are merged into these records, and records with a zero net change are removed. This avoids checking deletion requests again and modifying the base lists for every update.

To retrieve the current incoming neighbors of a destination, the index counts the occurrences of each source in the base list and adds the value in its delta record, if one exists. A source is returned if the resulting count is positive. For example, a base list $[u,u,w]$ and delta records $[(u,-2),(v,+1)]$ yield incoming sources $[v,w]$: both edges from $u$ have been removed, the edge from $w$ remains, and an edge from $v$ has been added. The base list remains unchanged during updates, while the delta records provide the current incoming view at query time.


\textbf{GPU descriptor and cache maintenance.} Reloading all GPU descriptors after a local update would transfer indices for many vertices whose adjacency remains unchanged. Source isolation makes these transfers unnecessary, while the source lists produced during adjacency maintenance identify the descriptors that need updating. Let $S^-$ and $S^+$ denote the sources changed during deletion and insertion. The system forms $S=S^-\cup S^+$ so that each source is included once, even if it changes in both stages. Sources with only unsuccessful deletion requests are excluded. After insertion, the CPU transfers one record per source in $S$, containing its final adjacency location, degree, and version, and the GPU writes these records into the corresponding descriptor entries. Sources outside $S$ retain valid descriptors because their adjacency has not moved or changed. Descriptor maintenance therefore follows the sources actually modified in the batch, including those updated in place or restored to their initial neighbors, without reloading the full array.

Each descriptor update occupies 24 bytes, giving a transfer volume of $24|S|$ bytes per batch. Descriptor traffic therefore scales with the number of changed sources rather than the total vertex count. Incoming-neighbor transfers, cached-list copies, and accesses to host adjacency incur separate costs.

The source list also limits the repair of cached adjacency. If a changed source has a valid GPU cache entry, the system copies its current neighbor list into the existing space when it fits. Otherwise, it attempts to place the list in unused space at the end of the cache and updates its cache index. If there is insufficient space, it invalidates the cached copy, allowing computation to read the current host adjacency. Each repair copies the changed source's entire neighbor list; cached lists of other sources are left in place during this step.

To prevent GPU computation from reading stale indices or cached neighbors, descriptor and cache updates run in one CUDA stream, and subsequent computation waits for them to finish. Replaced host blocks are reused only after earlier GPU readers have completed, preventing access to reclaimed adjacency storage.

Local repair can also avoid subsequent cache reorganization. After computation, the existing hotness-based policy selects the vertices to cache. If all selected vertices have valid cached adjacency and their number matches the previously published cache set, the system skips eviction, compaction, and loading. Otherwise, it performs the normal cache replacement procedure. This prevents a local topology change from unnecessarily triggering the full cache maintenance sequence, while still allowing replacement when the selected vertices change or local repair cannot retain a valid copy.

Together, source isolation, shared change records, and selective GPU updates confine adjacency relocation and cache repair to changed sources, while avoiding repeated processing of the same updates across topology indices. They reduce both data movement and the additional index maintenance.
